\subsection{Induced and coinduced root data}\label{sga3:XXI:6.5}

Let \(\cR=(M,M^\ast,R,R^\ast)\) be a root datum, and let
\(N\subset M\) be a subgroup containing the roots, that is,
\[
 \Gamma_0(R)\subset N\subset M.
\]
The canonical linear map \(i_N:N\to M\) gives, by transposition, a
linear map
\[
 {}^ti_N:M^\ast\longrightarrow N^\ast.
\]
Put
\[
 R_N=R,
 \qquad
 R_N^\ast={}^ti_N(R^\ast).
\]

\begin{lemma}\label{sga3:XXI:6.5.1}
The quadruple
\[
 \cR_N=(N,N^\ast,R_N,R_N^\ast)
\]
is a root datum, and \(i_N\) is a morphism.
\end{lemma}

\begin{proof}
First, \({}^ti_N\) induces an isomorphism from \(R^\ast\) onto
\(R_N^\ast\). Indeed, if \(\alpha,\beta\in R\) and
\[
 {}^ti_N(\alpha^\ast)={}^ti_N(\beta^\ast),
\]
then
\[
 (\alpha^\ast,x)=(\beta^\ast,x)
\]
for every \(x\in N\), in particular for \(x\in R\). Hence
\(\alpha=\beta\) by \ref{sga3:XXI:1.1.4}. The rest follows without
difficulty.
\end{proof}

\begin{definition}\label{sga3:XXI:6.5.2}
The root datum \(\cR_N\) is called the root datum
\emph{induced by \(\cR\) on \(N\)}.
\end{definition}

\begin{lemma}\label{sga3:XXI:6.5.3}
Let \(f:\cR'\to\cR\) be a morphism, and put
\[
 N=f(M')\subset M.
\]
Then \(f\) factors uniquely through \(i_N\).
\end{lemma}

In particular, up to isomorphism, the isogenies
\(\cR'\to\cR\) correspond bijectively to the finite-index subgroups of
\[
 M/\Gamma_0(R).
\]
This sharpens \ref{sga3:XXI:6.2.7}.

Now let \(N^\ast\subset M^\ast\) be a subgroup containing \(R^\ast\).
Define the root datum \emph{coinduced by \(\cR\) on \(N^\ast\)} by
\[
 \cR^{N^\ast}=(\cR^\ast_{N^\ast})^\ast.
\]
There is a canonical morphism
\[
 p^{N^\ast}:\cR\longrightarrow\cR^{N^\ast}.
\]

\begin{lemma}\label{sga3:XXI:6.5.4}
Let \(f:\cR'\to\cR\) be a morphism. There exist subgroups
\(N\subset M\) and \(N'^\ast\subset M'^\ast\) such that \(f\) factors
as
\begin{center}
\begin{tikzcd}[column sep=huge,row sep=large]
\cR' \arrow[r,"f"] \arrow[d,"p^{N'^\ast}"']
  & \cR \\
\cR'^{\,N'^\ast} \arrow[r,"f_0","\sim"']
  & \cR_N \arrow[u,"i_N"']
\end{tikzcd}
\end{center}
where \(f_0\) is an isomorphism.
\end{lemma}

\begin{proof}
Take \(N=f(M')\), as in \ref{sga3:XXI:6.5.3}. The resulting morphism
\(M'\to N\) is surjective, so its transpose is injective. Take its
image for \(N'^\ast\).
\end{proof}
